% related_work_cn.tex - 相关工作

\section{相关工作}

\subsection{微表情识别方法}

微表情识别方法可分为手工特征方法和深度学习方法。早期工作依赖时空描述符捕获微表情细微运动特征。LBP-TOP\cite{zhao2007lbptop}将经典LBP描述符扩展到三正交平面（XY、XT、YT），同时编码空间纹理和时间动态。基于光流的方法\cite{wang2015offapexnet}量化连续帧间像素级运动模式。这些方法计算高效，但受限于手工设计的特征表达能力。

深度学习时代带来了实质性改进。3D卷积网络\cite{tran2015c3d}建立了视频级表示学习基础。多尺度时序处理\cite{ben2022survey}在不同分辨率捕获动态。注意力机制具有重要影响：时空注意力网络在3D骨干中结合注意力模块。图注意力机制\cite{ben2022survey}建模面部肌肉运动模式。近期基于Transformer的方法\cite{liu2022videoSwin}利用自注意力实现层次化时空表示。

\subsection{双通路神经架构}

双通路设计在视频理解领域已有探索。双流网络\cite{simonyan2014twostream}分别处理空间（RGB）和时间（光流）信息后融合。SlowFast网络\cite{feichtenhofer2019slowfast}使用高分辨率慢通路和低分辨率快通路进行动作识别。在MER中，双流架构研究较少，大多数方法采用单骨干设计。

\subsection{混合专家模型}

混合专家框架\cite{jacobs1991moe}实现模块化学习，通过专门化子网络处理不同输入模式。稀疏门控\cite{fedus2022switch}通过仅激活相关专家实现高效扩展。在MER中，MoE具有吸引力，因为不同表情类别可能受益于专门化特征子空间——压抑微笑的面部动态与隐藏恐惧有质的差异。

\subsection{仿生计算}

仿生计算赋予人工系统生物神经处理的计算原理。双通路视觉假说启发了计算框架\cite{morris1999amygdala}。快皮层下路径（上丘-丘脑枕-杏仁核）快速检测情感显著刺激，慢皮层路径（V1-FFA-眶额皮层）提供精细分析。据我们所知，Censor是首个明确用不同架构偏置实现两条通路的MER系统。